\scanpage{187}
% archive foliation: 186
\uline{Proposition}\quad Soit $f:X-Y$ un morphisme, avec
\[
\left|
\begin{array}{ll}
1)&\deleted{f\text{ loc de présentation finie et loc quasi-fini}}\\[-.3ex]
 &\inserted{Y\text{ loc. noeth et univ. japonais}}\\
2)&\deleted{\unsure{Y}}\ X\quad\deleted{\text{irréductible}}\ 
\text{n'a qu'un nb fini de comp irred.}\\
3)&f\text{ loc. }\inserted{\uline{\text{de type fini}}}
\text{quasi-fini}\quad\deleted{\text{et, et}}\quad f\ \uline{\text{séparé}}
\end{array}
\right.
\]
\[
\left|\quad\text{Dans ces conditions, }f\text{ est \uline{de type fini}.}\right.
\]

On peut supposer $Y$ affine donc noeth, \unsure{\&}
$X$ intègre \deleted{et $f$ dominant}
\inserted{(car il suffit de voir que $f$ est quasi. cpt)}.
Soit $U$ un ouvert \inserted{affine} non vide de $X$, il est
quasi-fini séparé sur $Y$, donc par le main theorem il est
$Y$-isomorphe à un ouvert $\unsure{U|U'}$ d'un $Y$-\deleted{\unsure{entier alg}}
préschéma \uline{fini} $X'$, qu'on peut supposer intègre.

\margin{\diagram{\begin{tikzcd}[column sep=large,row sep=large]
&X''\arrow[dl,"\text{fini}"']\arrow[d]\\
X\arrow[d]&X'\arrow[dl,"\substack{\text{fini}\\\unsure{\text{surj}}}"]\\
Y&
\end{tikzcd}}}

Soit $X''$ le sous-préschéma fermé intègre de $X\times_YX'$
\deleted{\unsure{par prolongement}}
\inserted{adhérence de}
la ``diagonale'' de $U\times_YU'$, donc $X''$ s'envoie dans $X$ et $X'$.
Mais $X''-X$ est fini, et étant dominant, est \uline{surjectif}.
Il suffit donc de voir que $X''$ est de t.f. sur $Y$, ou même sur $X'$.
Or $X''-X'$ est un morphisme birationnel, avec $X''$ et $X'$ intègres.
Donc \deleted{\unsure{il}} y a un $X'$-morphisme
\deleted{\uline{surjectif}} $\widetilde X'-X''$, où $\unsure{X'|\widetilde X'}$
est le normalisé. Or $\unsure{X'|\widetilde X'}$ est fini sur $Y$,
on \unsure{gagne} \unsure{.....}
